\begin{abstract}
\ST{} feeds a reasoning model blends of token embeddings instead of single tokens, in the hope that the model can
carry several candidate continuations in superposition.
Recent work adds Gumbel noise to the blend weights to avoid greedy collapse.
With that noise, the top token of each blend is distributed exactly as a sampled token, so \GST{} is ordinary sampling
plus a residual weight on the other candidates.
We test superposition at the level of single tokens in DeepSeek-R1-Distill-Qwen-1.5B on GSM8K: does the residual % SETUP
survive the layers, persist in the cache, and change what the model does?
Exact 50/50 blends of two candidates, a forced superposition built at the input,
collapse gradually through the layers.
Information about the runner-up persists in the keys and values of layers 0 to 7 and shifts the projected % C21.10
next-token readout of later text toward it ($+0.44$, 95\% CI [0.35, 0.53], number vs number forks); % C21.7-8
on the model's own path the shift is not detectable.
In real runs the residual is small (median 0.12 at forks), and the top token changes as often as under % C22.12
sampling (0.45 vs 0.45). % C22.4
A coupled ablation with shared noise shows that the model reads the residual: it moves next-token predictions a
median 3.9 times as much as a random residual of the same size (95\% CI [3.1, 5.0]). % C23.15
Accuracy does not change detectably (88 vs 86 of 100; paired difference $+2$ points, 95\% CI [$-5$, $+9$]), % C23.1 C23.17
and the real and random residuals change sampling's final answer on largely the same problems % C23.18
(10 shared, 6 for each alone). % C23.18
Token-level superposition is therefore present in the input and read by the model, but we detect no use of it in
the answers of this model; whether whole reasoning paths are carried in superposition is not tested here.
The main predictions were committed before the runs, and every number from our experiments is recomputed from saved
outputs.
\end{abstract}
